\loai{Tính chiều dài thanh $\ell$ ở nhiệt độ t}
\begin{vidu}%[1P7-2.2K][Loai1] 
	Buổi sáng nhiệt độ là $15^\circ C$, chiều dài của thanh thép là 10 m. Hỏi buổi trưa khi nhiệt độ là $30^\circ C$ thì chiều dài của thanh thép là bao nhiêu? $\alpha = 1,1. 10^{-5}\ K^{-1}$.
	\choice
	{10 cm}
	{10,165 m}
	{\True 10,00165 m}
	{101,65 cm}
	\loigiai{
		Ta có chiều dài của thanh thép ở $30^\circ C$: $\ell=\ell_0.(1+\alpha.\Delta t)=10,00165\ m$.
	}
\end{vidu}
\loai{Tính độ nở dài $\Delta \ell$}
\begin{vidu}%[1P7-2.2K][Loai2] 
	Một thanh ray dài 10 m được lắp lên đường sắt ở nhiệt độ $20^\circ C$. Phải để hở một khe ở đầu thanh ray với bề rộng là bao nhiêu để khi nhiệt độ tăng lên $50^\circ C$ thì vẫn đủ chỗ cho thanh ray nở ra. Hệ số nở dài của chất làm thanh ray là $12.10^{-6}\ K^{-1}$.
	\choice
	{1,6 mm}
	{4,6 mm}
	{\True 3,6 mm}
	{2,6 mm}
	\loigiai{
		\begin{itemize}
			\item Khe hở là chung cho cả hai đầu thanh đối diện nhau nên khe hở phải đủ rộng để mỗi đầu nở ra $\dfrac{\Delta \ell}{2}$. Tức là khe hở phải có bề rộng $\Delta \ell$.
			\item Ta có: $\Delta \ell=\ell_0.\left(1+\alpha\Delta t\right)=10.12.10^{-6}.\left({50-20}\right)=3.10^{-3}\ m=3,6\ mm$
			\item Phải để hở đầu thanh ray 3,6 mm.
		\end{itemize}
	}
\end{vidu}

\loai{Tính nhiệt độ - nở dài}
\begin{vidu}%[1P7-2.2K][Loai4] 
	Mỗi thanh ray đường sắt ở nhiệt độ $15^\circ C$ có độ dài là 12,5 m. Nếu hai đầu các thanh ray chỉ đặt cách nhau 4,5 mm, thì các thanh ray này có thể chịu được nhiệt độ lớn nhất là bao nhiêu để chúng không bị uốn cong do tác dụng nở vì nhiệt? Cho biết hệ số nở dài của mỗi thanh ray là $12. 10^{-6}\ K^{-1}$.
	\choice
	{$25^\circ C$}
	{$20^\circ C$}
	{$35^\circ C$}
	{\True $45^\circ C$}
	\loigiai{
		Để thanh ray không bị cong khi nhiệt độ tăng thì độ tăng chiều dài của thanh phải bằng khoảng cách giữa hai đầu thanh ray.
		\begin{align*}
			\Delta \ell=\ell_0.\alpha.(t_2-t_1) \Rightarrow t_2=t_{\max}=\dfrac{\Delta \ell}{\alpha.\ell_1}+t_1 \Rightarrow t_{\max}=45^\circ.
		\end{align*}
	}
\end{vidu}

\loai{Tính thể tích vật ở nhiệt độ t}
\begin{vidu}%[1P7-2.2K][Loai5] 
	Một quả cầu bằng đồng thau có có đường kính 100 cm ở nhiệt độ $25^\circ C$. Tính thể tích của quả cầu ở nhiệt độ $60^\circ C$. Biết hệ số nở dài $\alpha =1,8.10^{-5}\ K^{-1}$.
	\choice
	{$0,325\ m^3$}
	{$0,425\ m^3$}
	{\True $0,525\ m^3$}
	{$0,625\ m^3$}
	\loigiai{
		\begin{itemize}
			\item Thể tích quả cầu ở $25^\circ C$: $V_0=\dfrac{4}{3}.\pi .R^3=\dfrac{4}{3}.3,14.\left({0,5}\right)^3=0,524\ m^3$.
			\item Mà $\beta =3\alpha =3.1,8.10^{-5}=5,4.10^{-5}\ K^{-1}$.
			\item Mặt khác 
			\begin{align*}
				\Delta V=\beta V_0\Delta t=5,4.10^{-5}.0,524.\left({60-25}\right) &\Rightarrow V_2=V_1+9,904.10^{-4}
				\Rightarrow V_2=0,5249904\ m^3.
			\end{align*}
		\end{itemize}
	}
\end{vidu}
\loai{Tính độ nở khối $\Delta V$}
\begin{vidu}%[1P7-2.2K][Loai6] 
	Tìm độ nở khối của một quả cầu nhôm bán kính 40 cm khi nó được đun nóng từ $0^\circ C$ đến $100^\circ C$, biết $\alpha =24.10^{-6}\ K^{-1}$.
	\choice
	{$0,93.10^{-3}\ m^3$}
	{$3,93.10^{-3}\ m^3$}
	{\True $1,93.10^{-3}\ m^3$}
	{$2,93.10^{-3}\ m^3$}
	\loigiai{
		\begin{itemize}
			\item Ta có thể tích của quả cầu ở $0^\circ C$: $V_0=\dfrac{4}{3}.\pi .R^3$.
			\item Độ nở khối của một quả cầu nhôm 
			\begin{center}
				$\Delta V=V-V_0=\beta V_0\Delta t=\dfrac{4}{3}.\pi .R^3.3.\alpha \Delta t$
				$\Rightarrow \Delta V=\dfrac{4}{3}.\pi .\left({0,4}\right)^3.3.24.10^{-6}.\left({100-0}\right)=1,93.10^{-3} \ m^3$.
			\end{center}
		\end{itemize}
	}
\end{vidu}
\loai{Tính khối lượng riêng}
\begin{vidu}%[1P7-2.2K][Loai7] 
	Vàng có khối lượng riêng là $1,93.10^4\ kg/m^3$ ở $30^\circ C$. Hệ số nở dài của vàng là $14,3.10^{-\text{ }6}\ K^{-1}$. Khối lượng riêng của vàng ở $110^\circ C$ là
	\motcot
	{$0,9234.10^4\ kg/m^3$}
	{$2,9234.10^4\ kg/m^3$}
	{\True $1,9234.10^4\ kg/m^3$}
	{$3,9234.10^4\ kg/m^3$}
	\loigiai{
		 Ta có công thức tính khối lượng
			\begin{align*}
				m=D_0.V_0=D .V\Rightarrow D =\dfrac{{V_0}}{V}.D_0=\dfrac{{D_0}}{{1+\beta .\Delta t}}
				\Rightarrow D =\dfrac{{1,{93.10}^4}}{{1+3.14,{3.10}^{-6}.\left({110-30}\right)}}=1,9234.10^4\ kg/m^3.
			\end{align*}	
\begin{note}
	Khối lượng của chất luôn không đổi.
\end{note}
	}
\end{vidu}
\loai{Bài toán tổng hợp}
\begin{vidu}%[1P7-2.2G][Loai8] 
	Một lá nhôm hình chữ nhật có kích thước $2\ m\ x\ 1\ m$ ở $0\doC$. Đốt nóng tấm nhôm tới $400\doC$ thì diện tích tấm nhôm sẽ là bao nhiêu? $\alpha =25.10^{-6}\ K^{-1}$.
	\choice
	{$0,04\ m^2$}
	{$1,04\ m^2$}
	{\True $2,04\ m^2$}
	{$3,04\ m^2$}
	\loigiai{
		\Binhluan{
		\begin{itemize}
			\item Ta có diện tích $S=a.b$ 
			\item Chiều dài của cạnh thứ nhất thay đổi theo nhiệt độ: 
			\begin{align*}
				a=a_0(1+\alpha \Delta t)=2.\left({1+{25.10}^{-6}.\left({400-0}\right)}\right)=2,02\ m.
			\end{align*}
			\item Chiều dài của cạnh thứ nhất thay đổi theo nhiệt độ: 
			\begin{align*}
				b=b_0(1+\alpha \Delta t)=1.\left({1+{25.10}^{-6}.\left({400-0}\right)}\right)=1,01\ m.
			\end{align*}
			\item Khi đó: $S=a.b=2,02.1,01=2,04\ m^2$.
		\end{itemize}
	}
{Ta cũng có thể dùng công thức tính hệ số nở diện tích: $\gamma=2\alpha$:
\begin{align*}
	S=S_0.\left(1+\gamma\Delta t\right)
\end{align*}
}
	}
\end{vidu}